\subsection{B.1. Categories and the four local descriptions}

For the mixed-characteristic comparison statements, $K$ is a complete
discretely valued field of characteristic zero, with perfect residue
field $\kappa$ of characteristic $p$ and valuation ring $\mathcal O_K$.
Put $K_0=W(\kappa)[1/p]$, $C=\widehat{\overline K}$ and
$G_K=\Gal(\overline K/K)$. The absolute ramification
$e=[K:K_0]$ is any finite positive integer. All modern statements below
include $p=2$; restrictions of the older sources are stated separately.
The residue field need not be finite. Denote by
$\operatorname{BT}(\mathcal O_K)$ and $\operatorname{BT}(K)$ the
categories of finite-height $p$-divisible groups; a subscript $\Q$
means that isogenies have been inverted.

We use the covariant lattice
$T_pG=\varprojlim_nG[p^n](\overline K)$ and
$V_pG=T_pG\otimes\Q_p$. The \emph{Hodge--Tate exponent} $r$ means
the exponent in $C(r)$; thus $\Q_p(1)$ has exponent $+1$.
For a group over $\mathcal O_K$, $V_pG$ has exponents $0,1$.
For its dual representation $V_pG^\vee$, the de Rham filtration has
jumps $0,1$. This dualization specifies the variance connecting the
two conventions.

The page supplies a single Hodge--Tate grading. Turning it into a
bigrading $(r,w-r)$ requires a separately specified pure weight $w$;
a general $p$-divisible group comes with no such specified motivic
weight. Likewise ``pseudo-motifs'' on 18:5 is not accompanied by a
category or morphisms. The following exact classifications concern
the named categories, and cannot constitute an essential-image
classification for an undefined category of pseudo-motives.

\begin{result}{B.1.1}{The generic-fibre Tate functor}{Proved}
There are equivalences
\[
 \operatorname{BT}(K)\simeq
       \Rep_{\Z_p}^{\mathrm{free,cont}}(G_K),\qquad
 \operatorname{BT}(K)_{\Q}\simeq
       \Rep_{\Q_p}^{\mathrm{fd,cont}}(G_K).
\]
Thus all continuous representations occur on the generic fibre,
including ones that never become Hodge--Tate after finite extension.
\end{result}
\begin{proof}
Finite group schemes over a characteristic-zero field are \emph{\'etale}.
Finite \emph{\'etale} descent identifies $G[p^n]$ with its finite
$G_K$-module of geometric points. For a height-$h$ group these are
free $\Z/p^n$-modules and their compatible maps recover $T_pG$.
Conversely a free continuous lattice $T$ gives the finite \emph{\'etale}
groups attached to $T/p^nT$, with the usual inclusions and
multiplication maps. After the faithfully flat extension
$\overline K/K$, the divisibility axioms are those of
$(\Q_p/\Z_p)^h$. These constructions and their Hom maps are
inverse. This is also the explicit descent description in
Tate~\cite[\S1.4 and \S2.1]{Tate1967}.

Every continuous finite-dimensional representation has a stable
lattice. Its image is compact; the translates of a fixed lattice
are bounded between two fixed scalar multiples of it. Their sum
is therefore a finitely generated stable lattice. Rational maps
become lattice maps after multiplication by a power of $p$, proving
the rational equivalence.

For an explicit exception to potential Hodge--Tate-ness, take
$K=\Q_p(\mu_p)$ for odd $p$, and $K=\Q_2(i)$ for $p=2$.
Let $\chi$ be the cyclotomic character, set $s=1/(p+1)$, and define
\[
                  \rho_s(g)=\exp(s\log\chi(g)).
\]
The image of $\chi$ is contained in $1+p\Z_p$ for odd $p$ and
$1+4\Z_2$ for $p=2$, so these series converge and give a continuous
unit-valued character. It defines a height-one \emph{\'etale}
$p$-divisible group over $K$. Exactly $\rho_s^{p+1}=\chi$.
If $\rho_s$ became Hodge--Tate over a finite $L/K$, its $C$-line
would have an integer exponent $r$, so its $(p+1)$st tensor power
would have exponent $(p+1)r$. The cyclotomic exponent is $1$.
Tate's Theorem~2 in \S3 gives $H^0(G_L,C(m))=0$ for $m\ne0$,
so different exponents cannot be isomorphic. It follows that
$(p+1)r=1$, a contradiction. The argument holds for every finite
$L$.
\end{proof}

\begin{result}{B.1.2}{Good and potentially good reduction; all primes}{Proved}
The functor
\[
 \operatorname{BT}(\mathcal O_K)\longrightarrow
                         \Rep_{\Z_p}^{\mathrm{free,cont}}(G_K)
\]
is fully faithful. Its essential image consists of every free
stable lattice in a crystalline representation with Hodge--Tate
exponents in $\{0,1\}$. After inverting isogenies its image is exactly
that class of crystalline representations. The potentially good
generic-fibre category has image exactly the representations
becoming crystalline with these exponents over a finite extension.
\end{result}
\begin{proof}
Tate~\cite[\S4.2, Theorem 4 and Corollary 1]{Tate1967} gives
\[
 \operatorname{Hom}_{\mathcal O_K}(G,H)
       \xrightarrow{\sim}\operatorname{Hom}_{G_K}(T_pG,T_pH).
\]
His normal-noetherian-domain hypothesis includes this DVR. Tensor
with $\Q_p$ to obtain the rational Hom assertion. The crystalline
essential-image assertion follows from the height-one classification
in Kim~\cite[Theorem 2.3 and Theorem 4.1]{Kim2012}; the former
records Kisin's rational result and the integral result for $p>2$,
and the latter supplies the integral $p=2$ case. The filtration
condition is translated explicitly in B.1.5 below.

For clarity, rational existence also supplies every stable lattice.
Start with one group and one lattice $T_0$ in the desired
representation. Up to scaling any other stable lattice is a finite
index superlattice of $T_0$. The corresponding finite subgroup of
the generic fibre has finite flat schematic closure in a suitable
$G_0[p^n]$; flatness follows from torsion-freeness over the DVR.
The quotient by that finite flat subgroup is a $p$-divisible group
with the prescribed superlattice. Scaling back is an isomorphism
of the corresponding rational representation. Equivalently apply
the integral height-one classification directly.

If a generic-fibre group acquires an $\mathcal O_L$-model, its
representation restricts to the indicated crystalline one over $L$.
Conversely use B.1.1 to construct the $K$-group from the given
$G_K$-representation, and use the good-reduction classification
over $L$ to obtain its model. This proves the asserted image.
It does not produce an $\mathcal O_K$-model. Descent data cannot
be omitted: over $K=\Q_5$, the trivial character and the
nontrivial quadratic character cut out by $\Q_5(\sqrt5)$ become
equal after that extension, but their Hom space as $G_K$-modules
is zero. Thus the restricted representation without descent
does not classify the original $K$-object.
\end{proof}

\begin{result}{B.1.3}{The different-prime and equal-characteristic readings}{Refuted; exact replacements proved}
For $\ell=p$, description $5'$ is the functor in B.1.1--B.1.2.
For $\ell\ne p$, the Tate functor of a $p$-divisible group is zero,
has essential image only zero, and is not faithful. In equal
characteristic $p$, even the $p$-adic geometric Tate functor is
not faithful on all $p$-divisible groups.
\end{result}
\begin{proof}
On a $p^n$-torsion group, integers $a,b$ with $a\ell+bp^n=1$
show that multiplication by $\ell$ is invertible. Hence
$G[\ell^m]=0$ and $T_\ell G=0$. The nonzero identity of
$\Q_p/\Z_p$ is killed. If the input is an abelian variety,
$A[\ell^\infty]$ is instead its $\ell$-primary summand; the
primary-torsion decomposition gives the exact relation to the
$p$-primary input.

For the equal-characteristic reading take $K=\mathbf F_p((t))$
and $G=\mu_{p^\infty}$. Its geometric $p$-power torsion points
are all the identity, so $T_pG=0$, while $\id_G\ne0$.
Dieudonn\'e or crystalline data, including the connected part,
provide the characteristic-$p$ replacement discussed in B.3.4.
\end{proof}

\begin{result}{B.1.4}{Descriptions $2'$ and $3'$ and their Hom spaces}{Proved}
For a Hodge--Tate representation set
\[
 D_{\mathrm{HT},r}(V)=(V\otimes_{\Q_p}C(-r))^{G_K},\qquad
 \bigoplus_rD_{\mathrm{HT},r}(V)\otimes_K C(r)
       \xrightarrow{\sim}V\otimes_{\Q_p}C.
 \tag{B.1.4a}
\]
The map is evaluation. If the continuous $G_K$-action is retained,
adjoining this canonical grading gives a fully faithful functor
on $\operatorname{BT}(\mathcal O_K)_\Q$, with essential image
the crystalline exponent-$\{0,1\}$ representations equipped with
(B.1.4a). If only the underlying $\Q_p$-space and its $C$-grading
are retained, the functor is faithful but not full, even with good
reduction. The analogous statements hold with ``potentially
crystalline'' on the potentially good domain.
\end{result}
\begin{proof}
Tate~\cite[\S4.1, Theorem 3 and Corollary 2]{Tate1967} supplies
the decomposition for $p$-divisible groups; dualizing his
$\operatorname{Hom}(T(G),C)$ formulation gives (B.1.4a) with our
exponents. Equivariant maps preserve its invariant spaces, so
the first assertion follows from B.1.2.

For failure of fullness take $\kappa=\mathbf F_p$ and two rank-one
unramified characters with arithmetic Frobenius eigenvalues $1$
and $a$, where $a=1+p$ for odd $p$ and $a=5$ for $p=2$.
The character $\widehat\Z\to\Z_p\to\Z_p^\times$ is continuous.
Finite \emph{\'etale} descent over $\mathcal O_K$ gives their
height-one \emph{\'etale} $p$-divisible groups. Both have only
exponent zero, so their bare graded objects are the same line
$\Q_p$ with its whole $C$-extension in degree zero. Its bare Hom
space is $\Q_p$. Equivariance requires $(a-1)x=0$, so the Galois
Hom space is zero. The forgetful functor is nevertheless faithful
because its underlying vector-space map determines a morphism.

There is a complete explicit parametrization of the bare essential
image. Range over the weakly admissible filtered $\varphi$-modules
$D$ with $\Fil^0D_K=D_K$ and $\Fil^2D_K=0$, form
$V_{\mathrm{cris}}(D)^\vee$ by B.2.1, form (B.1.4a), and forget
the Galois action. Up to graded $\Q_p$-isomorphism these are
exactly the objects in the bare image, by B.1.5. This is a
parametrization by filtered Frobenius linear algebra and specified
period rings, without referring to an unknown group. For the
potentially good image range over finite $L/K$ with the required
$G_K$-descent until the final forgetful step.

Finally Hodge--Tate-ness descends across finite extensions. Replace
one by a finite Galois extension. The invariant spaces defining
each exponent are stable under the finite quotient and are
semilinear spaces over $L$; finite Galois descent preserves their
dimensions. Their descended sum is (B.1.4a) over $K$.
The character of B.1.1 therefore proves that the phrase ``after
finite extension'' does not cover arbitrary generic-fibre groups.
\end{proof}

\begin{result}{B.1.5}{Description $6'$: the exact admissible category}{Proved}
Let $\operatorname{MF}_K^\varphi$ consist of finite-dimensional
$K_0$-spaces $D$ with bijective $\sigma$-semilinear $\varphi$, and
an exhaustive separated decreasing filtration on $D_K=D\otimes_{K_0}K$.
Morphisms commute with $\varphi$ and preserve the filtration. Put
\[
 t_H(D)=\sum_i i\dim_K\operatorname{gr}^iD_K,\qquad
 t_N(D)=v_p(\det\varphi).
\]
Weak admissibility means $t_H(D)=t_N(D)$ and
$t_H(D')\leq t_N(D')$ for every $\varphi$-stable subspace $D'$
with its induced filtration. There is an anti-equivalence
\[
 \operatorname{BT}(\mathcal O_K)_\Q^{\mathrm{op}}\simeq
 \{D\in\operatorname{MF}_K^{\varphi,\mathrm{wa}}:
                          \Fil^0D_K=D_K,\ \Fil^2D_K=0\},
 \quad G\longmapsto D_{\mathrm{cris}}(V_pG^\vee).
 \tag{B.1.5a}
\]
The $K_0$-space and its scalar extension to $K$ are essential parts
of this statement; a Frobenius lift on $K$ need not exist.
\end{result}
\begin{proof}
Changing a basis multiplies the determinant of the semilinear
$\varphi$ by $\sigma(\det A)/\det A$, of valuation zero, so $t_N$
is well-defined. Colmez--Fontaine~\cite[Theorem A, Proposition 4.2,
Theorem 4.3 and its corollary]{ColmezFontaine2000} proves weak
admissibility equals admissibility. Restrict its filtered
$(\varphi,N)$ equivalence to $N=0$ to get crystalline
representations. The indicated dual converts the exponent
condition into filtration jumps $0,1$. B.1.2 gives precisely the
$p$-divisible groups; the two fully faithful functors give all
Hom identifications in (B.1.5a).

The raw filtered data on the page do not impose admissibility.
For example $D=K_0e$, $\varphi(e)=e$, $\Fil^1D_K=D_K$,
$\Fil^2D_K=0$ has $t_H=1$, $t_N=0$, and cannot occur.
For failure of a Frobenius lift to $K$, take
\[
 p=2,\quad K_0=\Q_2(\zeta_3),\quad
 K=K_0(\pi),\qquad \pi^3=2\zeta_3.
\]
The polynomial is Eisenstein; $e=3$ and the residue field is the
perfect field $\mathbf F_4$. The Witt Frobenius sends
$\zeta_3\mapsto\zeta_3^2$. An extension to $K$ would make
$z=\sigma(\pi)/\pi$ a unit satisfying $z^3=\zeta_3$.
Reduction is impossible, since every nonzero element of
$\mathbf F_4$ has cube $1$, while $\overline\zeta_3\ne1$.
Thus the scalar-extension morphism $D\to D_K$ is the precise
replacement for placing both structures on the same $K$-space.
\end{proof}

For the integral classification, choose a uniformizer $\pi$ with
Eisenstein polynomial $E(u)$ and put
$\mathfrak S=W(\kappa)[[u]]$, $\varphi(u)=u^p$. A height-one
Kisin module is finite free over $\mathfrak S$, with injective
linearized Frobenius whose cokernel is killed by $E(u)$.
Breuil's Theorems 1.2 and 4.2.2.9 give his strongly divisible
classification for $p>2$; his Theorem 1.4 also assumes
$\kappa\subset\overline{\mathbf F}_p$ for the displayed
crystalline essential-image result~\cite{Breuil2000}.
Kim's Theorems 2.3 and 4.1 give the all-prime statement over the
perfect residue field used here. Proposition 4.2 identifies the
associated Dieudonn\'e module after evaluation on the PD-envelope.
No Fontaine--Laffaille bound $e<p-1$ is imposed. At $p=2$ the
naive map to the Breuil $A_{\mathrm{cris}}$ realization can fail
to be an isomorphism, as Kim notes immediately before Theorem 4.1;
the corrected equivalence is required. The original Kisin 2006
source was not obtained in a readable full format in this session;
the opened Kim statements support the asserted all-prime form.

\subsection{B.2. Replacements for \texorpdfstring{$1'$ and $4'$}{1-prime and 4-prime}}

\begin{result}{B.2.1}{The period-ring pair and reconstruction}{Proved}
Put $A_{\inf}=W(\mathcal O_C^\flat)$, with
$\theta:A_{\inf}\to\mathcal O_C$;
$B_{\dR}^+$ is the $\ker\theta$-adic completion of $A_{\inf}[1/p]$,
and $B_{\dR}=B_{\dR}^+[1/t]$. Here
$t=\log[\epsilon]$ for compatible roots of unity,
$g(t)=\chi(g)t$ and $\Fil^iB_{\dR}=t^iB_{\dR}^+$.
Let $A_{\mathrm{cris}}$ be the $p$-completed divided-power envelope
of $\ker\theta$, compatible with the standard divided powers of
$p$, and set $B_{\mathrm{cris}}=A_{\mathrm{cris}}[1/p,1/t]$.
The ring $B_{\mathrm{st}}$ also has monodromy $N$ satisfying
$N\varphi=p\varphi N$; its embedding into $B_{\dR}$ uses the
chosen logarithm/uniformizer convention.

For $B=B_{\dR},B_{\mathrm{cris}},B_{\mathrm{st}}$, put
$D_B(V)=(B\otimes_{\Q_p}V)^{G_K}$, over $K$ in the de Rham case
and $K_0$ in the other two. On the $B$-admissible category,
defined by equality of this dimension with $\dim_{\Q_p}V$,
evaluation is the comparison isomorphism
\[
           B\otimes D_B(V)\xrightarrow{\sim}B\otimes_{\Q_p}V.
 \tag{B.2.1a}
\]
It respects the filtration, Frobenius and monodromy whenever
those structures occur. The inverse for an admissible filtered
$(\varphi,N)$-module is
\[
 V_{\mathrm{st}}(D)=
  (B_{\mathrm{st}}\otimes_{K_0}D)^{\varphi=1,N=0}
  \cap\Fil^0(B_{\dR}\otimes_KD_K),
 \tag{B.2.1b}
\]
inside their common $B_{\dR}$ extension; use $B_{\mathrm{cris}}$
and omit $N$ for the crystalline inverse.
\end{result}
\begin{proof}
The inverse identities and full faithfulness are
Colmez--Fontaine, Proposition 4.2 and Theorem 4.3, with its
crystalline corollary. Thus (B.2.1a)--(B.2.1b) give the exact
functors and scalar extension replacing the proposed local
``double viseur''. The underlying filtered space alone loses
information through the explicit functor
$(D,\varphi,\Fil)\mapsto(D_K,\Fil)$: the two unramified
characters in B.1.4 have the same one-dimensional filtered de
Rham space, but distinct crystalline Frobenius and zero Galois
Hom space. Retaining $V$ and the comparison, or retaining the
admissible Frobenius data for the inverse, restores the stated
reconstruction.
\end{proof}

\begin{result}{B.2.2}{Geometric and relative comparisons}{Proved}
For $X/K$ proper smooth, all $p$, all finite $e$, perfect
$\kappa$ and no reduction assumption, there is a filtered
$G_K$-equivariant comparison
\[
 H^i_{\mathrm{et}}(X_{\overline K},\Q_p)\otimes_{\Q_p}B_{\dR}
  \simeq H^i_{\dR}(X/K)\otimes_KB_{\dR}.
 \tag{B.2.2a}
\]
A proper smooth $\mathcal O_K$-model upgrades it to the
$B_{\mathrm{cris}}$ comparison with the special fibre's rational
crystalline cohomology. A proper flat regular semistable model
with reduced normal-crossing special fibre gives the
$B_{\mathrm{st}}$ comparison with Hyodo--Kato cohomology,
Frobenius and monodromy. Potential semistability is obtained
after a finite extension; descent information must still be kept.

For the relative form let $f:X\to Y$ be proper smooth between
smooth rigid-analytic varieties over $K$, and let $\mathbb L$ be
a de Rham lisse $\Z_p$-sheaf on $X$ with associated filtered
connection $(\mathcal E,\nabla,\Fil)$. Assume
$R^if_*\mathbb L$ is lisse, as holds for the analytification of
the algebraic situation. It is de Rham, with associated filtered
connection
\[
                         R^if_{\dR*}(\mathcal E,\nabla,\Fil).
 \tag{B.2.2b}
\]
For the constant sheaf this is the Gauss--Manin connection and
Hodge filtration on relative de Rham cohomology.
\end{result}
\begin{proof}
Scholze~\cite[Theorem 1.6 and Corollary 1.8]{Scholze2013}
gives (B.2.2a) for proper smooth rigid spaces; analytification
gives the algebraic assertion. Tsuji's opened author survey
\cite[Theorem 1.1 and Appendix Theorem A1]{Tsuji2002Survey}
states the semistable theorem and its de Rham consequence.
The good-reduction case has $N=0$ and its crystalline comparison.
The potential statement is also explicit in
\cite[Theorem 2.2.1 and Corollary 2.2.3]{Tsuji1998ICM}.

Equation (B.2.2b) is the application of Scholze's Theorem 1.10
to the stated sheaf and morphism. The comparison uses the
relative pro-\emph{\'etale} sheaf $\mathcal O\mathbb B_{\dR}$,
whose canonical connection corresponds to the displayed
Gauss--Manin connection; the filtration satisfies Griffiths
transversality. This proves the indicated local replacement
of $4'$ after scalar extension to that period sheaf.
The three-page 2016 corrigendum~\cite{Scholze2016Corrigendum}
was checked: its corrected coverings and the proper
$p$-completion in $\mathcal O\mathbb B_{\dR}^+$ are used;
it retains the comparison results.
\end{proof}

\begin{result}{B.2.3}{Integral lattices, torsion and the prismatic classification}{Proved within the stated scope}
For a proper smooth formal $\mathfrak X/\mathcal O_C$, integral
comparison is carried by the perfect Frobenius complex
$R\Gamma_{A_{\inf}}(\mathfrak X)$. Its derived specializations
recover crystalline cohomology over $W(\kappa_C)$, de Rham
cohomology over $\mathcal O_C$, crystalline cohomology over
$A_{\mathrm{cris}}$, and the generic fibre's $\Z_p$-\emph{\'etale}
cohomology after inverting $\mu=[\epsilon]-1$.
If the special fibre's crystalline cohomology in degrees $i$
and $i+1$ is $p$-torsion-free, integral crystalline $H^i$, with
Frobenius, is reconstructed from the \emph{\'etale} lattice
and its $B_{\dR}^+$-lattice furnished by the comparison.
For the base $\operatorname{Spf}(\mathcal O_K)$, prismatic
$F$-crystals in finite locally free modules classify exactly
free $\Z_p$-lattices in crystalline representations, for all
$p$, arbitrary finite $e$ and perfect $\kappa$.
\end{result}
\begin{proof}
The four derived isomorphisms are BMS
\cite[Theorem 1.8]{BMS2018}. Its Theorems 1.4, 1.6 and 1.7
give the precise reconstruction: the pair $(T,\Xi)$, with
$\Xi\subset T\otimes B_{\dR}$ a $B_{\dR}^+$-lattice, gives
a finite free Breuil--Kisin--Fargues module, whose crystalline
specialization is the asserted $H^i$. The two adjacent
torsion hypotheses are retained.

The final equivalence is Bhatt--Scholze
\cite[Theorem 1.2, equivalently Theorem 5.6]{BhattScholzePrismaticCrystals}.
Evaluation on the Breuil--Kisin prism
$(W(\kappa)[[u]],(E(u)))$ gives the integral Frobenius-module
data; the crystal condition supplies the compatible changes of
thickening. These assertions give both the structures and their
specified comparison functors.
\end{proof}

The integral object in this result is a complex, and its
specializations are derived tensor products. Replacing them by
ordinary tensor products of individual cohomology groups
requires additional flatness or torsion hypotheses, since the
derived base-change spectral sequence otherwise has Tor terms.
The filtered complex of differentials of a chosen smooth model
gives an integral Hodge filtration, but its cohomology may have
torsion and that filtration need not consist of direct summands.
The rational relative theorem gives the filtered vector bundle
and integrable connection of B.2.2. The local prismatic
classification just stated does not assert that all integral
cohomology groups of all relative smooth proper families are
finite locally free crystals.

\subsection{B.3. The four questions of 18:43}

\begin{result}{B.3.1}{Problem 1a: the two meanings of lattices}{Proved and literal unscaled equality refuted}
For $k\hookrightarrow\C$, the complex comparison is
\[
 H^i_{\dR}(X/k)\otimes_k\C
       \simeq H^i_{\B}(X(\C),\Q)\otimes_{\Q}\C.
\]
If ``r\'eseaux'' means the two field forms in this common space,
their position is precisely the period point of A. If it means
integral lattices, the de Rham lattice needs an integral model,
and even the standard model of $\mathbf P^1$ does not give
unscaled equality: in the conventions of A and F,
\[
                     \alpha(e_{\B})=(2\pi i)^{-1}e_{\dR}.
\]
Thus the transported Betti lattice is
$(2\pi i)^{-1}\Z e_{\dR}$. Tensor powers give the corresponding
Tate factors in every degree.
\end{result}
\begin{proof}
The comparison is integration of algebraic forms, and the explicit
codimension-one cycle calculation in F.1 gives its value on the
chosen generators. Hence the asserted comparison of lattices
follows by applying that linear map to each integer multiple;
it also supplies the required counterexample to unscaled equality.
No equality of two integral lattices has been inferred from the
complex scalar-extension isomorphism alone.
\end{proof}

\begin{result}{B.3.2}{Problem 1c: real structures and the splitting}{Proved}
For a real model, with a rational Betti basis and real de Rham
basis, the rational cohomological involution $f_\infty$ satisfies
\[
 \kappa_D\alpha=\alpha f_\infty\kappa_V,
 \qquad \overline\alpha=\alpha f_\infty.
\]
Consequently the de Rham real subspace is
\[
 D_{\R}=\alpha\bigl((V^+\otimes\R)\oplus
                                      i(V^-\otimes\R)\bigr),
\]
whereas transport of the Betti real form is $\alpha(V\otimes\R)$.
The Hodge splitting is obtained from $F_D$ and
$\alpha f_\infty\alpha^{-1}F_D$ by the opposed-filtration
construction of A.2.
\end{result}
\begin{proof}
The conjugation comparison of A.2.1, now using the real descent
identification of the conjugate object, gives the displayed
semilinear identity. Writing $v=v_++v_-$, the condition
$f_\infty\overline v=v$ means $v_+$ has real coefficients and
$v_-$ has purely imaginary coefficients; this is exactly the
displayed real subspace. The opposite flag and its normalized
splitting are the maps already proved in A.2.5. For a motive
without real descent, conjugation instead has the typed target
$c^*M$ of A.2.1; it does not supply an involution of the same
object. Algebraicity of the resulting splitting remains the
explicit period question of A.4(c), and is not implied by these
real-structure equations.
\end{proof}

\begin{result}{B.3.3}{Problem 3: the precise local comparison inputs}{Proved within the specified readings}
For proper smooth $X/K$, reading ``$p$-adic cohomology'' as
$H^i_{\mathrm{et}}(X_{\overline K},\Q_p)$ gives (B.2.2a).
Reading it as special-fibre crystalline cohomology requires a
good-reduction model and gives the $B_{\mathrm{cris}}$
comparison. A semistable model replaces that special-fibre
object by log crystalline/Hyodo--Kato cohomology, adds $N$, and
uses $B_{\mathrm{st}}$.
\end{result}
\begin{proof}
Apply the three cases of B.2.2. The Hyodo--Kato map to de Rham
cohomology is Tsuji~\cite[Theorem 3.2.1]{Tsuji1998ICM}; its
uniformizer convention is part of the map. The full filtered
$(\varphi,N)$-module and its inverse (B.2.1b) record the
representation despite that convention. This specifies the
extra input and the coefficient-extension morphism for each
reading. No assertion for arbitrary unspecified spaces or raw
unscaled integral lattices follows from the phrase alone.
\end{proof}

\begin{result}{B.3.4}{Problem 4: characteristic-$p$ cohomology and the requested structures}{Proved within each stated scope}
Over perfect $\kappa$, with $K_0=W(\kappa)[1/p]$, the following
table supplies the characteristic-$p$ replacements. Every
Frobenius is semilinear for the indicated base Frobenius.
\end{result}
\begin{center}\small
\begin{longtable}{p{.17\textwidth}p{.24\textwidth}p{.20\textwidth}p{.28\textwidth}}
\toprule
Theory and scope&Integral or rational object&Frobenius&Filtration and connection\\
\midrule
\endhead
Crystalline, $X/\kappa$ smooth proper&Perfect $W(\kappa)$-complex;
finite cohomology possibly with torsion&Isogeny after inverting
$p$&A smooth lift gives its de Rham filtration; relative crystals
give the integrable connection\\
Rigid, $X/\kappa$ separated of finite type&Finite-dimensional
$K_0$-cohomology, also with compact supports; an integral lattice
is extra data&Bijective after linearization&Overconvergent
coefficient isocrystals carry connections; Hodge filtration
needs additional lift or geometric data\\
Prismatic, $X$ smooth $p$-adic formal over $A/I$, $(A,I)$ a bounded
prism&Derived integral $A$-complex, perfect for proper $X$&Linearized
Frobenius becomes invertible after inverting $I$&Nygaard
filtration and de Rham specialization; crystal descent supplies
compatible changes of thickening\\
\bottomrule
\end{longtable}
\end{center}
\begin{proof}
For the prismatic row use Bhatt--Scholze
\cite[Theorem 1.8(1)--(6), Theorem 1.16 (=15.3),
Example 1.9(1)]{BhattScholzePrisms}. At the crystalline prism
$(W(\kappa),(p))$, their crystalline comparison identifies
crystalline cohomology with the Frobenius pullback of the
prismatic complex. Since the Witt Frobenius is an automorphism,
inverse scalar twist undoes this pullback. This is the exact
morphism relating the two integral theories. Perfectness over
the DVR gives finite crystalline cohomology, and the Frobenius
isogeny gives invertibility after inverting $p$. The de Rham
specialization and the filtered complex of a lift give the
indicated Hodge filtration. A Nygaard filtration is a specified
additional filtration on a prismatic construction; it has not
been identified with a chosen lift's Hodge filtration without
the comparison morphism.

For rigid cohomology, Kedlaya
\cite[Theorems 1.2.1--1.2.2]{Kedlaya2006} gives finite-dimensional
ordinary and compactly supported cohomology, even with
overconvergent $F$-isocrystal coefficients. For constant
coefficients Frobenius invertibility is stated and proved in
the opened author text of D\'eglise
\cite[Proposition 6.2 and Remark 6.3]{Deglise2005}: smooth
proper-hypercover descent reduces to smooth schemes, and
compactification and localization handle compact supports.
Berthelot~\cite[Proposition 1.9]{Berthelot1997} identifies the
smooth proper rigid theory with crystalline cohomology after
inverting $p$. Thus the table records an explicit scalar
extension in their common scope.

The prismatic \emph{\'etale} comparison in Theorem 1.8(4)
requires a perfect prism and refers to the characteristic-zero
generic fibre. A characteristic-$p$ scheme at the crystalline
prism has empty such generic fibre; that statement is not an
identification of all crystalline cohomology with the special
fibre's $p$-adic \emph{\'etale} cohomology. This qualification
prevents a false extension of the mixed-characteristic
comparison to a different input.
\end{proof}

\subsection{B.4. The two splittings on 18:6}

\begin{result}{B.4.1}{Compatibility and the exact cocycle}{Proved; implication from centralization alone refuted}
Let $1\to H\to E\to\pi\to1$ be a topological group extension,
with $H=G'(C)$ and continuous sections $\sigma_1,\sigma_2$.
Use the stated comparison $\alpha$ to transport
$u:\pi\to G(\Q_p)$ into $H$, and write
$j:\Q_p^\times\to H$ for $j_2$ and
$\chi:\pi\to\Q_p^\times$ for the character. The page's
centralizations imply that $\sigma_1u$ and $\sigma_2j\chi$
are sections. Their equality is an additional compatibility.
Under that compatibility,
\[
 \varphi(\gamma)=\sigma_1(\gamma)^{-1}\sigma_2(\gamma)
                    =u(\gamma)j(\chi(\gamma)^{-1}),
 \tag{B.4.1a}
\]
and the exact right-cocycle law is
\[
 \varphi(\gamma\delta)=
 \sigma_1(\delta)^{-1}\varphi(\gamma)\sigma_1(\delta)
                                                 \varphi(\delta).
 \tag{B.4.1b}
\]
Its identity fibre is the equalizer of the sections, a subgroup
that need not be normal.
\end{result}
\begin{proof}
The commutations assumed on the page allow the two factors in
each modified section to multiply separately, proving their
section laws. If those sections agree, cancel $\sigma_1(\gamma)$
on the left and $j(\chi(\gamma))$ on the right to get
(B.4.1a); commutation of $u$ with $j$ is not needed.
Substitute $\sigma_2(\gamma)=\sigma_1(\gamma)\varphi(\gamma)$
in the two original section laws to get (B.4.1b). Two
homomorphisms' equalizer is closed under products and inverses,
proving the subgroup assertion.

For failure of the extra compatibility, take $\pi=C_2$ written
additively, $H=C^\times$, $E=H\times\pi$,
\[
 \sigma_1(\gamma)=(1,\gamma),\quad
 \sigma_2(\gamma)=((-1)^\gamma,\gamma),\quad
                  u=1,\quad\chi=1,\quad j=\id.
\]
All centralizations and section laws hold, but the modified
sections differ at $\gamma=1$. This checks every hypothesis
except the missing equality and proves it cannot be deduced.
\end{proof}

\begin{result}{B.4.2}{The kernel identity and descent of $\chi$}{Proved; unqualified claim refuted}
Under (B.4.1a), the unconditional kernel identity is
\[
             \ker\chi\cap\ker\varphi=\ker\chi\cap\ker u.
\]
It equals $\ker u$ exactly when $\ker u\subset\ker\chi$.
The character $\chi$ descends to $\Pi_1=\pi/\ker u$ exactly
under that same condition. A sufficient typed motivic condition
is a character $\varepsilon:G\to\Gm$ with
$\chi=\varepsilon\circ u$.
\end{result}
\begin{proof}
When $\chi(\gamma)=1$, formula (B.4.1a) says
$\varphi(\gamma)=u(\gamma)$, proving the first equality.
Its left side is contained in $\ker\chi$, which proves
necessity of the inclusion; that inclusion also proves
sufficiency. A homomorphism factors through $\pi/\ker u$
exactly when it kills $\ker u$, and then its formula is
$\chi_1([\gamma])=\chi(\gamma)$.

Even compatibility does not imply the inclusion. In the same
direct product as B.4.1, set
\[
 \chi(\gamma)=(-1)^\gamma,\quad u=1,\quad
 \sigma_1(\gamma)=(1,\gamma),\quad
 \sigma_2(\gamma)=(\chi(\gamma)^{-1},\gamma),\quad j=\id.
\]
Then $\sigma_1u=\sigma_2j\chi$, but
$\ker u=C_2$ and $\ker\chi\cap\ker\varphi=\{0\}$.
The construction also works with $\pi=G_K$ and its cyclotomic
character. The finite example and the right-cocycle identities
were checked by exact arithmetic in the accompanying program.
\end{proof}

\begin{result}{B.4.3}{Analyticity and the precise relation to Sen}{Proved within the named scope}
If $\pi$ is profinite, $u$ continuous, and $G$ a
finite-dimensional linear algebraic $\Q_p$-group, then
$\Pi_1=\pi/\ker u$, with its quotient topology, is a
$p$-adic analytic group for every $p$. A Hodge--Tate grading
gives $h:\mathbb G_{m,C}\to\GL(V_C)$ and its differential
$\Theta=dh(1)$, acting by the exponent on each graded summand.
The group identities above do not by themselves prove Sen's
minimal-algebraic-group assertion.
\end{result}
\begin{proof}
The image $u(\pi)$ is compact, hence closed in $G(\Q_p)$.
The continuous bijection $\pi/\ker u\to u(\pi)$ is a
homeomorphism by compactness and Hausdorffness. The closed
subgroup theorem for $p$-adic Lie groups, in the opened author
proof \cite[Corollary 19.14]{Pstragowski2023}, then supplies its
analytic structure; this argument requires no Hodge--Tate
hypothesis. The formula for $\Theta$ is differentiation of
the actual grading cocharacter.

The primary treatment opened for the comparison with Sen is
Serre~\cite[\S1.1, \S1.3 Theorems 1 and $1'$, \S1.4
Theorem 2 and its equivalent forms]{Serre1979HodgeTate}.
Its residue field is explicitly algebraically closed. In that
scope the Galois image of a Hodge--Tate representation is
open in an algebraic group, and its connected algebraic
envelope is the smallest algebraic $\Q_p$-group containing
the grading cocharacter after $C$-extension. Serre attributes
the result to Sen 1973, \S6. More generally the Sen
infinitesimal statement concerns inertia, not an unrestricted
unramified direction. Neither an equality of modified
sections nor their kernel identity proves that minimality.
The following explicit example also disproves replacing
inertia by the full Galois image without the residue-field
hypothesis.

Let $K=\Q_5$ and take the unique root
$s=\sqrt6\in\Z_5$ with $s\equiv1\pmod5$.
Hensel's lemma applies because $2$ is a unit, and $s\notin\Q$
since $6$ is not a rational square. Define an unramified
two-dimensional representation by sending arithmetic Frobenius
to $\operatorname{diag}(6,\exp(s\log6))$. Its full compact
image is
\[
 z\longmapsto
     \operatorname{diag}(\exp(z\log6),\exp(sz\log6)),
                       \qquad z\in\Z_5.
\]
It has Lie algebra $\Q_5\operatorname{diag}(1,s)$, whereas
both characters have Hodge--Tate exponent zero, so $h=1$
and $\Theta=0$. Both arise from \emph{\'etale} $p$-divisible
groups over $\Z_5$, which verifies the claimed Hodge--Tate
property. Inertia is trivial, in agreement with the inertia
statement.

Moreover the full-image Lie algebra is not algebraic. A
connected one-dimensional subgroup of $\Gm^2$ is killed by
a nonzero integral character $x^my^n$. Its tangent line
would satisfy $m+ns=0$, impossible for irrational $s$.
If an algebraic subgroup of $\GL_2$ had this Lie algebra,
its identity component would be a torus in the diagonal
centralizer of any nonzero semisimple element of the line,
so the same contradiction applies. The image's Zariski
closure is the whole $\Gm^2$: a proper connected subtorus
would impose the same impossible relation. Thus the
full-image formulation is actually false, rather than merely
unproved by the section calculation.
\end{proof}

\begin{result}{B.4.4}{An Artin--Tate proalgebraic counterexample with both splittings}{Refuted; exact counterexample proved}
The unrestricted proalgebraic version of the analyticity
assertion in I.P33 is false, even for an inverse system of
good-reduction Artin--Tate realizations over $\Q_p$. It remains
false after imposing all the compatibility and character-descent
identities of B.4.1--B.4.2.
\end{result}
\begin{proof}
Fix a prime $\ell\ne p$. Let $K_m/\Q_p$ be the unramified
extension of degree $\ell^m$, and let
$r_\ell:G_{\Q_p}\twoheadrightarrow\Z_\ell$ be the associated
quotient. The motives $h^0(\operatorname{Spec}K_m)$ have the
regular representations of $\Z/\ell^m$; their integral models
are finite \emph{\'etale} over $\Z_p$. Adjoin the Tate line
with character $\chi_p$. Write
\[
 A=\varprojlim_m(\Z/\ell^m)_{\Q_p},\qquad
 G=A\times\Gm,\qquad
 u(\gamma)=(r_\ell(\gamma),\chi_p(\gamma)).
 \tag{B.4.4a}
\]
The first factor is a proconstant affine group scheme; it has
$A(F)=\Z_\ell$ for every extension field $F/\Q_p$. Equip this
factor with its inverse-limit topology. The finite-dimensional
Artin representations are unramified and crystalline of
exponent zero; the Tate factor is crystalline of exponent one.
Thus every finite tensor construction in this category is
crystalline and has the indicated good-reduction origin.

Every $\Q_p(\mu_{p^n})/\Q_p$ is totally ramified: the
polynomial $\Phi_{p^n}(1+T)$ is Eisenstein, including $p=2$
when the extension is nontrivial. A finite unramified and a
totally ramified extension have trivial intersection. Taking
their composita at each finite level and passing to inverse
limits therefore proves the exact image formula
\[
                  u(G_{\Q_p})=\Z_\ell\times\Z_p^\times.
 \tag{B.4.4b}
\]
At level $m$ this image is Zariski dense in
$(\Z/\ell^m)_{\Q_p}\times\Gm$: each component contains the
infinite subgroup $\Z_p^\times$ of $\Gm$. Consequently $G$
is precisely the proalgebraic group of these representations,
rather than merely a larger target containing their image.

Here are all the group data in B.4.1. Set $C=\widehat{\overline{\Q}_p}$,
$H=G(C)=\Z_\ell\times C^\times$, and use the natural action
$\gamma(a,z)=(a,\gamma z)$. Form the semidirect extension
\[
 1\longrightarrow H\longrightarrow H\rtimes G_{\Q_p}
                  \longrightarrow G_{\Q_p}\longrightarrow1.
\]
Take $G'=G$; the comparison of these commutative group schemes
is the identity, since conjugation by any comparison point
acts trivially on a commutative group. Define
\[
\begin{aligned}
 j(t)&=(0,t),&\varepsilon(a,z)&=z,\\
 \sigma_1(\gamma)&=((0,1),\gamma),&
 \sigma_2(\gamma)&=((r_\ell(\gamma),1),\gamma).
\end{aligned}
 \tag{B.4.4c}
\]
All maps are continuous and have the stated domains:
$j:\Q_p^\times\to H$ is induced by a cocharacter, and
$\varepsilon:G\to\Gm$ is a character. The constant first
factor makes $\sigma_2$ a homomorphic section. The action of
$G_{\Q_p}$ fixes $u(G_{\Q_p})$ and $j(\Q_p^\times)$ pointwise;
$H$ is abelian. These facts prove both centralizations, and
direct multiplication gives
\[
 \sigma_1u=\sigma_2j\chi_p,\qquad
 \varphi(\gamma)=(r_\ell(\gamma),1),\qquad
 \chi_p=\varepsilon\circ u,\qquad
 \ker u=\ker\chi_p\cap\ker\varphi.
\]
Thus even the necessary character-descent condition holds.
Compactness identifies the quotient $\Pi_1=G_{\Q_p}/\ker u$
homeomorphically with (B.4.4b).

A compact $p$-adic analytic group has an open pro-$p$ subgroup;
see \cite[Theorem 8.32 and Corollary 8.34, pp.200--201]{AuditDDMS1999}.
Every continuous homomorphism from a pro-$p$ group to
$\Z_\ell$ is zero, because its finite images are simultaneously
$p$-groups and $\ell$-groups. Hence any pro-$p$ subgroup of
$\Z_\ell\times\Z_p^\times$ lies in
$\{0\}\times\Z_p^\times$, which has empty interior.
No such subgroup is open. This proves that $\Pi_1$ has no
$p$-adic analytic structure compatible with its quotient
topology. At every finite Artin level the image is instead
$(\Z/\ell^m)\times\Z_p^\times$, a $p$-adic analytic group,
so the failure is exactly the passage to this infinite inverse
system. All assertions hold also for $p=2$.
\end{proof}

The exact proved content of 18:6 is therefore the cocycle
identity under the stated compatibility, the kernel equality
under its necessary and sufficient additional inclusion,
and analyticity under continuity and finite dimensionality.
The grading differentiates to the operator occurring in Sen
theory. Lead H3 interpreted as a proof of the minimality
theorem from those few identities is refuted by the missing
conditions and the full-image counterexample. Interpreted as
a historical anticipation of Sen's theorem, it remains a
priority claim requiring documentary work; the algebra is
not a proof of that historical claim.

\subsection{B.5. Source and computation obligations retained}

The modern positive results above have opened full sources
and exact theorem locators. The prompt also requires opening
particular historical originals. Fontaine--Messing 1987,
Faltings 1989, Tsuji's 1999 Inventiones article, Kisin 2006,
and Sen's 1973 Annals paper were not obtained in readable
full text during the attempts recorded here. Their metadata
or publisher previews do not meet that requirement. The
opened replacements are respectively the stated comparison
theorems of Scholze, Tsuji's own full 1998 and 2002 accounts,
Kim's full classification paper, and Serre's primary
Hodge--Tate treatment. The historical source obligations
therefore remain unresolved in this incomplete record;
no full-text opening of those originals is claimed.

The exact check script
\texttt{comparison\_counterexamples.py} records all inputs
and outputs for the two $C_2$ section examples, their
cocycle and kernel identities, the $\mathbf F_4$ cube
obstruction, and the weak-admissibility arithmetic.
It also constructs the Hensel root of $x^2-6$ through
modulus $5^{12}$ as an illustration. That finite precision
is not used to prove irrationality or nonalgebraicity;
those are the exact arguments in B.4.3.
